\chapter{From the Nine Muses and Four Elements to RealQM and AI}\label{ch:muses}
\index{Empedocles}\index{Muses}\index{Love and Strife}\index{density functional theory!Hohenberg--Kohn theorem}
\index{saving the phenomena}\index{Kepler, J.}\index{artificial intelligence}

\begin{quote}
\itshape Hear first the four roots of all things \ldots\ now through Love all coming together into one, now
again each borne apart by the hatred of Strife.\\
\normalfont\hfill --- after Empedocles, fragments B6 and B17
\end{quote}
\bigskip

The Greeks sought to understand the world through a few fundamental constituents interacting through a few
universal principles. Empedocles named four elements and two forces, Love and Strife. This chapter follows
that programme through the history of physics to RealQM, in which electrons are charge densities in
three-dimensional space interacting through Coulomb's law. It places density functional theory, the working
method of present-day quantum chemistry, in the same history, and asks whether AI, by taking over the
solution of equations, may return physics to its first question: what is a physically meaningful model of
reality?

\section{The Muses and the elements}\label{sec:muses-elements}

The ancient Greeks approached both the human mind and the physical world through a small number of
fundamental concepts. The nine Muses represented the forms of intellectual and artistic activity: Calliope
of epic poetry, Clio of history, Erato of love poetry, Euterpe of music, Melpomene of tragedy, Polyhymnia of
sacred song, Terpsichore of dance, Thalia of comedy---and Urania of astronomy~\cite{Hesiod}. That astronomy
had its own Muse says something important: science was part of human culture, not a separate technical
enterprise, and understanding nature was bound up with the perception of beauty, harmony and order.

\section{Empedocles: four roots and two forces}\label{sec:empedocles}

Around 450\,BC Empedocles of Akragas in Sicily proposed that everything is composed of four
\emph{roots}---the word \emph{element} came later---namely earth, water, air and fire. Their mixing and
separation are governed by two opposing principles:
\begin{itemize}
\item \textbf{Love} (\emph{Philotes}), which brings things together;
\item \textbf{Strife} (\emph{Neikos}), which drives them apart~\cite{DK51,KRS83,Inwood01}.
\end{itemize}
This is a remarkably modern form of physical theory: a few fundamental constituents, interacting through
universal principles, generate the immense diversity of the observable world. Empedocles wrote his theory as
a poem and invoked the Muse Calliope to assist him~\cite{DK51}. Physics, philosophy and poetry were still one
enterprise.

\section{Love and Strife as gravitation and Coulomb interaction}\label{sec:love-strife}

In the physical world envisaged by RealQM two interactions take centre stage:
\begin{itemize}
\item \textbf{gravitation}: attraction between masses;
\item \textbf{Coulomb interaction}: attraction between opposite charges, repulsion between equal ones.
\end{itemize}
Gravitation is Love in its purest classical form: it gathers matter into planets, stars and galaxies.
Coulomb interaction carries both principles. The attraction between positive nuclei and negative electrons
brings matter together; the repulsion between charges of equal sign counteracts unlimited concentration and
organises matter in space. Together with wave dynamics, this interplay of attraction and repulsion produces
the structures of atoms, molecules and condensed matter.

Both laws are inverse-square laws~\cite{Newton1687,Coulomb1785}, and both can be written as Poisson
equations in three-dimensional space,
\begin{equation*}
\Delta\phi = 4\pi G\rho_{\mathrm{mass}}, \qquad -\Delta V = 4\pi\rho_{\mathrm{charge}},
\end{equation*}
so that across 2{,}500 years there is a conceptual line from Love and Strife to gravitation and Coulomb
interaction. It is not a literal identification, but a shared vision of nature organised by a few universal
principles of interaction.

\section{Strife has two faces}\label{sec:strife-two-faces}
\index{Pauli exclusion!as a wall cost}\index{helium dimer}

Recent computation adds a twist that Empedocles might have enjoyed. Coulomb repulsion is not the whole of
Strife.

Take two helium atoms. Each is a closed shell: two electrons of opposite spin forming a tight pair. In
nature two helium atoms brought close repel each other; there is no He$_2$ molecule in the ordinary sense.
In the basic form of RealQM, where electron domains meet at free walls, the computation instead finds a
false bond of about half a hartree at a separation of two bohr. Coulomb repulsion between the electron
clouds is present, but it is not enough. This is the missing repulsion between electron groups recorded in
Chapter~\ref{ch:limits}.

What is missing is the other face of Strife: \emph{exclusion}~\cite{Pauli1925}. Two electrons of the same
spin cannot occupy the same region of space. In RealQM this is not an abstract symmetry of a wave function
but a spatial fact: electrons of the same spin meet at a wall, and the wall must carry a cost---the energy
of the part of the electron's density that the wall cuts off. Outside its own region a bound electron's
density would continue as a tail decaying at the rate $\kappa=\sqrt{-2\varepsilon}$ set by its own
eigenvalue $\varepsilon$; the energy of that tail, charged at the wall, is $\tfrac{\kappa}{2}\int u^2\,dS$,
a Robin condition $\partial_n u+\kappa u=0$. With that cost, set by the electron's own energy and nothing
fitted, the false bond disappears and helium repels as it should. The same holds for two beryllium atoms,
and for two hydrogen atoms with parallel spins, whose repulsive triplet curve the wall reproduces at short range (too repulsive at long range, where the cost decays as $e^{-\kappa R}$ and the true repulsion as $e^{-2\kappa R}$). Between
the two electrons of a bonding pair, which have opposite spins and may share space, the wall stays free.

So in RealQM Love is Coulomb attraction, and Strife comes in two forms: Coulomb repulsion between charges,
and exclusion between electrons of the same spin. Both are statements about matter in three-dimensional
space.

\section{From Empedocles to modern physics}\label{sec:empedocles-modern}

The history of physics can be read as a sequence of attempts to carry out the programme of Empedocles:
\begin{itemize}
\item Democritus: atoms and the void;
\item Newton: matter interacting through gravitation;
\item Maxwell: electric and magnetic fields~\cite{Maxwell1865};
\item Einstein: space, time, matter and energy;
\item Schr\"odinger: matter described by wave mechanics~\cite{Schrodinger1926};
\item the Standard Model: elementary particles and quantum fields.
\end{itemize}
Modern physics recognises four fundamental interactions: gravitation, electromagnetism, and the strong and
weak nuclear interactions. RealQM, with its extensions to the nucleus (Chapter~\ref{ch:nucleus}) and to
cosmology (Chapter~\ref{ch:cosmology}), explores a more economical programme built on gravitation and
Coulomb interaction. Whether such a reduction can account for the full range of observed phenomena is an
open scientific question.

\section{Standard quantum mechanics: simplicity lost?}\label{sec:simplicity-lost}

Standard quantum mechanics (StdQM), originating with Schr\"odinger in 1926, gives a formally compact
description of atoms and molecules. But the wave function of $N$ electrons depends on $3N$ spatial
coordinates, not on the three coordinates of physical space. Its direct numerical representation requires
resources growing exponentially with the number of electrons---Kohn's ``exponential wall''~\cite{Kohn1999}---and
a century after its introduction the physical meaning of a function on $3N$-dimensional configuration space
remains disputed.

This is a striking contrast to the programme of Empedocles: a formally simple equation produces a
representation of nature of enormous computational complexity.

Density functional theory (DFT)~\cite{HK64,KS65} brings the computation back to three dimensions, and it is
today the workhorse of chemistry and materials science. But it does so by collecting everything that cannot
be computed into one unknown term, the exchange--correlation functional, which is defined through the
$3N$-dimensional wave function and is in practice a chosen formula. DFT's own leading developers note that
there is no systematic way to improve these formulas, and that some fitted to energies have drifted away from
the physics~\cite{Burke2012,Medvedev2017}. DFT gives good numbers; it does not claim that its objects are
physical. It is a three-dimensional computational method for a theory posed in $3N$ dimensions.

\section{Density functional theory: saving the phenomena}\label{sec:saving-phenomena}

\paragraph{The theorem.} DFT is presented as resting on the Hohenberg--Kohn theorem~\cite{HK64}: the
ground-state energy is a functional of the electron density alone, so that DFT is ``exact in principle''.
The argument is a one-line contradiction. Suppose two different external potentials gave the same
ground-state density; apply the variational principle to each ground state and add the two inequalities,
and one obtains $E_1+E_2<E_1+E_2$. Later work by Levy and Lieb~\cite{Levy1979,Lieb1983} repaired its
technical assumptions, and as mathematics it stands.

But the argument uses only two things: that the ground state minimises the energy, and that the external
potential enters linearly, as $\int v\rho$. It uses nothing of the kinetic energy, of the Coulomb interaction
between electrons, of the number of electrons, or of three-dimensional space. It therefore holds unchanged
for \emph{any} interaction between the particles---Coulomb, any other law, or none. A statement that holds
whatever the interaction cannot say anything about the effects of this one. It is an existence statement
reached by contradiction: it shows that a functional exists and says nothing about what it is or how to
compute it. Everything DFT must actually get right---how electrons avoid each other, exchange, correlation,
binding---lies in the exchange--correlation functional, about which the theorem is silent. That functional
is defined through the $3N$-dimensional Schr\"odinger equation and cannot be computed without solving it.
An exact theory that cannot be computed makes no prediction and cannot be tested; what is tested is always
one of the chosen formulas.

\paragraph{What the theorem does.} Its roles are four, and only the first is pure prestige. It gave the
density approach a mathematical warrant, and with it respectability in physics: no longer a crude model in
the line of Thomas and Fermi, but ``exact in principle''. It gave permission to seek functionals of the
density. It supplies a target: exact properties of the defined functional (its scaling, its limits) are used
to constrain approximations such as PBE~\cite{PBE1996}. And it underlies the Kohn--Sham
construction~\cite{KS65}, whose fictitious non-interacting electrons reproduce the density and carry the
kinetic energy. The physics DFT does contain---kinetic energy, the attraction of the nuclei, the classical
Coulomb energy of the density---is taken over from the Schr\"odinger equation; the many-electron physics that
remains is supplied by formulas, chosen by model systems, by imposed constraints, or by fitting to molecular
data~\cite{Becke1993}.

History bears this out. Chemists adopted DFT broadly in the 1990s, some thirty years after the theorem, when
gradient-corrected and hybrid functionals began to give useful accuracy. Adoption followed the numbers, not
the theorem. Kohn's Nobel Prize of 1998 was in chemistry, shared with Pople for computational methods: DFT
was honoured as a method, not as physics. The usual presentation has the order reversed: the theorem is
offered as the foundation and the approximations as an unfortunate necessity, whereas the approximations
carry all the content and the theorem supplies reassurance.

\paragraph{A classical analogue.} The fair classical analogue is not dogma: a scripture is accepted on
authority and cannot be refuted, while DFT's formulas are tested against measurement every day and their
failures are published by their own developers. The analogue is Greek astronomy and its explicit doctrine of
\emph{saving the phenomena}~\cite{Duhem1908}. An astronomical model did not need to be physically true; it
had to reproduce the observed positions of the planets. Ptolemy's circles were calculating devices, not the
mechanism of the heavens, and every new discrepancy was met by adjusting the devices. DFT practises the same
astronomy. Its Kohn--Sham orbitals are declared fictitious, as the circles were; its functional is judged by
whether it reproduces energies and geometries; new failures are met by new functionals, of which some two
hundred have been assessed side by side~\cite{MHG17}; and a mathematical warrant, the Hohenberg--Kohn
theorem, supplies the legitimacy that the doctrine of perfect circles once gave.

A modern analogue, from fluid mechanics, has the same structure. Averaging the Navier--Stokes equations
leaves an unknown term, the Reynolds stress, exactly defined by the full equations and uncomputable without
solving them; it is modelled by closures with fitted constants, of which there is likewise a zoo. DFT's
exchange--correlation term is the Reynolds stress of quantum chemistry.

In 1609 Kepler called his \emph{Astronomia Nova} a \emph{physica coelestis}, a celestial
physics~\cite{Kepler1609}: he sought the causes of the planetary motions, not only their positions. His
ellipses were at first no more convenient than the established tables, but they were physics, and Newton
could build on them. RealQM takes Kepler's side: a model offered as physical, accepting worse numbers for the
present in exchange for errors that point to causes. The analogy concerns method, not truth. Ptolemy's
cosmos was false, whereas the Schr\"odinger equation behind DFT is not known to be wrong: DFT saves the
phenomena of a correct but uncomputable theory. But saving the phenomena is what it does.

\section{RealQM: recovering the simplicity of nature}\label{sec:muses-realqm}

RealQM is an attempt to recover the original ambition of natural philosophy with the means of modern
computational mathematics. In RealQM:
\begin{itemize}
\item electrons are charge densities distributed in three-dimensional physical space;
\item electrons occupy non-overlapping domains with free, moving boundaries;
\item electrons and nuclei interact through Coulomb attraction and repulsion, each electron with the others
but not with itself;
\item atomic and molecular dynamics are governed by time-dependent partial differential equations in physical
space.
\end{itemize}
Instead of one wave function in $3N$ dimensions, an $N$-electron system is represented by $N$ interacting
three-dimensional densities. The computational work can then grow essentially linearly with the number of
electrons, given suitable algorithms.

Every object in RealQM is offered as physical, and every approximation is a statement about how electrons
behave. That makes the model testable piece by piece. When the basic model bound two helium atoms, the
failure pointed to one specific claim---that a wall between electrons costs nothing---and the correction was
a physical rule. This is the difference between a model of physics, whose wrong numbers point to the physics
that is missing, and a model for numbers, which can be refitted indefinitely without teaching anything about
nature. The first is science in the sense of Empedocles; the second, however useful, is closer to Ptolemy's
epicycles.

RealQM's present numbers are honest rather than final. For simple molecules its bond lengths are within a
few per cent of experiment, at the level of the simplest density functionals, while its binding energies are
still off by tens of per cent (Chapters~\ref{ch:hydrides} and~\ref{ch:limits}). The validity of the
programme must be settled by computation and comparison with observation.

In the broader picture Coulomb interaction organises matter on the small scale and gravitation on the large.
From molecules to galaxies, nature would be understood as matter in physical space, interacting through two
laws.

\section{AI and the revival of natural philosophy}\label{sec:muses-ai}

AI adds a new dimension to this story. Mathematical modelling of the world has three steps:
\begin{enumerate}
\item formulating a question about the world;
\item expressing it in a mathematical model;
\item solving the resulting equations.
\end{enumerate}
AI is rapidly transforming the third step and increasingly assists with the second. This may shift the
attention of physics back to the first and most fundamental question: \emph{what is a physically meaningful
mathematical model of reality?}

The challenge is then not only to solve the equations of StdQM more efficiently, but to ask whether other
equations give a simpler, more transparent and more computable description of the physical world. AI makes it
possible to explore and test such alternatives at a pace that was unthinkable a few years ago: the work on
the wall cost reported above---testing a hypothesis, finding where it fails, locating the missing physics and
testing the correction---was carried out in a dialogue between a human and an AI
(Chapter~\ref{ch:collab}).

\section{Back to the nine Muses}\label{sec:back-to-muses}

The nine Muses remind us that understanding involves not only logic and calculation but imagination,
intuition, beauty and creativity. Empedocles reminds us that the ambition of physics is to explain the complex
through the simple. Newton and Coulomb gave attraction and repulsion precise mathematical form; RealQM seeks to
build on these laws through computational models in physical space; and AI may now provide the means to
explore that vision on an unprecedented scale.

Could the era of AI bring a renaissance of natural philosophy, uniting physical intuition, mathematical
modelling and computational exploration? Could RealQM be a step in that direction?

After 2{,}500 years the vision of Empedocles remains remarkably attractive: the immense complexity of the
world emerging from a few constituents interacting through Love and Strife.
